\subsection{Middleware de comunicación}

El middleware multi-protocolo de SensorWave implementa un puente de traducción transparente entre CoAP, HTTP, MQTT y NATS, permitiendo interoperabilidad entre dispositivos IoT heterogéneos. Mediante un patrón de broadcast selectivo, los mensajes publicados en cualquier protocolo se redistribuyen automáticamente a suscriptores en todos los demás protocolos, manteniendo un espacio de nombres unificado de tópicos. La arquitectura utiliza una interfaz común para clientes (\texttt{Publicar}, \texttt{Suscribir}, \texttt{Desuscribir}) y servidores específicos por protocolo que normalizan mensajes entrantes. Las implementaciones incluyen Server-Sent Events (SSE) para HTTP, observación CoAP (RFC 7641), integración con brokers MQTT y NATS. El sistema adopta consistencia eventual con propagación asíncrona, priorizando latencia baja sobre ordenamiento estricto, apropiado para telemetría IoT donde datos recientes invalidan valores previos. Esta capa de interoperabilidad permite que dispositivos con recursos limitados utilicen CoAP mientras que nodos edge y despachadores exponen APIs HTTP/REST, desacoplando las decisiones de protocolo en diferentes niveles del sistema sin sacrificar interoperabilidad global.

\subsection{Nodos edge}

Una vez que el middleware normaliza los mensajes entrantes, los nodos edge son responsables del almacenamiento, procesamiento y gestión de las series temporales. El componente edge constituye el elemento fundamental de la arquitectura de edge computing del sistema, implementando un gestor completo de series temporales con capacidades de almacenamiento local, procesamiento de reglas en tiempo real y sincronización con la nube.

Cada nodo edge tiene un identificador único, una base de datos local embebida PebbleDB para persistencia local, y opcionalmente una conexión a FoundationDB para sincronización con la infraestructura en la nube. El diseño dual de almacenamiento permite al nodo operar tanto en modo autónomo (solo con PebbleDB) como en modo conectado, donde puede sincronizar metadatos y migrar datos históricos a la nube según sea necesario.

La gestión de series temporales implementa un sistema de buffering multinivel con compresión sofisticada. Cada serie temporal mantiene su propio buffer en memoria, donde las mediciones entrantes se acumulan hasta alcanzar un tamaño de bloque configurable. Este diseño permite optimizar el throughput de escritura mediante operaciones por lotes. El sistema de compresión opera en dos niveles: primero, se aplican algoritmos específicos para marcas temporales (DeltaDelta) y valores (configurable entre DeltaDelta, RLE, empaquetado de bits, diccionario, o codificación categórica); segundo, el bloque resultante se comprime mediante algoritmos de propósito general como LZ4, Zstandard, Snappy o Gzip, permitiendo adaptar el balance entre velocidad y tasa de compresión según los requisitos de cada serie.

El componente edge incorpora un motor de reglas integrado que permite la evaluación de condiciones y ejecución de acciones directamente en el nodo edge, reduciendo la latencia de respuesta crítica. El motor soporta condiciones complejas sobre series individuales o grupos de series, con operadores de comparación estándar y funciones de agregación temporal (promedio, máximo, mínimo, suma, conteo). Las condiciones pueden especificarse mediante patrones de path con wildcards o filtros por tags, permitiendo reglas que evalúen dinámicamente conjuntos de series sin necesidad de enumerarlas explícitamente. La lógica booleana configurable (AND/OR) entre condiciones permite expresar reglas complejas de múltiples variables. Las acciones ejecutables incluyen logging, envío de alertas y activación de actuadores, con un sistema de ejecutores extensible que permite registrar nuevos tipos de acciones en tiempo de ejecución.

El subsistema de consultas implementa operaciones de lectura eficientes sobre los datos comprimidos. Las consultas por rango temporal utilizan un sistema de salto anticipado de bloques (\textit{early skipping}) que examina los metadatos de las claves de los bloques antes de descomprimirlos, evitando procesamiento innecesario de bloques que no intersectan con el rango temporal solicitado. Las consultas de punto único priorizan el buffer en memoria antes de acceder al almacenamiento persistente, optimizando el caso común de consultas sobre datos recientes.

La comunicación con la nube se realiza mediante dos mecanismos principales: (1) sincronización periódica de metadatos de series hacia FoundationDB, ejecutada cada vez que se crea una nueva serie, y (2) capacidad de migración de datos históricos a la nube que transfiere bloques completos de datos comprimidos desde PebbleDB a FoundationDB y los elimina localmente, liberando espacio en dispositivos edge con almacenamiento limitado.

\subsection{Nodos despachadores}

Para coordinar el acceso a datos distribuidos entre múltiples nodos edge, el sistema incorpora despachadores que actúan como puntos de entrada unificados para consultas. El despachador funciona como un componente stateless de enrutamiento de consultas en la arquitectura distribuida. Su diseño minimalista refleja una filosofía de separación de responsabilidades donde el despachador no mantiene datos de series temporales, sino que actúa exclusivamente como directorio dinámico y proxy de consultas.

Cada nodo despachador mantiene un mapa en memoria de nodos edge registrados, sincronizado periódicamente con FoundationDB. Esta arquitectura de consistencia eventual permite que el despachador descubra automáticamente nodos edge que se registran en FoundationDB sin necesidad de configuración estática o registro explícito. Cada entrada de nodo incluye su identificador único, dirección IP, puerto HTTP y el mapa completo de series que hospeda, información suficiente para enrutar consultas.

El mecanismo de enrutamiento de consultas sigue un patrón de petición-respuesta asíncrono. Cuando el despachador recibe una consulta para una serie específica, primero localiza el nodo edge que hospeda esa serie. Una vez identificado el nodo, el despachador establece una comunicación HTTP temporal con el nodo edge, publica la solicitud de consulta en un tópico específico formateado como \texttt{consulta/<nodoID>/<tipoConsulta>}, y se suscribe a un tópico de respuesta único generado con UUID (\texttt{respuesta/<IDPeticion>}). Este diseño basado en tópicos evita colisiones entre consultas concurrentes y permite timeouts configurables.

El despachador soporta tres tipos de consultas: consultas por rango temporal, que retornan todas las mediciones entre dos marcas temporales; consulta del último punto, optimizada para monitoreo en tiempo real; y consulta del primer punto, útil para determinar el inicio del historial de datos. Cada tipo de consulta utiliza un parser específico para deserializar la respuesta del nodo edge, proporcionando una interfaz fuertemente tipada a los clientes del despachador.

El diseño stateless del despachador ofrece ventajas significativas en escalabilidad horizontal: múltiples instancias de despachador pueden ejecutarse simultáneamente sin necesidad de coordinación entre ellas, ya que toda la información de estado compartida reside en FoundationDB. Esta arquitectura permite balanceo de carga trivial y tolerancia a fallos mediante simple replicación activa-activa del componente despachador.

\subsection{Interacción entre edge y despachador}

La arquitectura completa implementa un patrón de descubrimiento de servicios donde los nodos edge operan de manera autónoma, publicando sus capacidades (series disponibles) en FoundationDB, mientras que los despachadores consumen esta información para construir dinámicamente su vista del sistema distribuido. Esta separación permite que los nodos edge se concentren en las operaciones intensivas de I/O y procesamiento local, mientras que los despachadores manejan la complejidad de consultas distribuidas y agregación de resultados de múltiples nodos, estableciendo una clara división de responsabilidades que facilita la escalabilidad independiente de cada capa del sistema.
